\begin{abstract}
Lewy body dementia (LBD)---encompassing dementia with Lewy bodies (DLB) and Parkinson's disease dementia (PDD)---is the second most common cause of neurodegenerative dementia, yet it remains profoundly underdiagnosed and entirely lacking in disease-modifying treatments. The molecular complexity of LBD, driven by convergent dysregulation across the genome, epigenome, transcriptome, proteome, and metabolome, demands an integrated multi-omics strategy that no single-layer approach can provide. This proposal outlines a computational research project designed to integrate publicly available LBD multi-omics datasets in order to discover molecular biomarkers and elucidate pathological mechanisms. Building on recent landmark studies---including a 2025 genome-wide association study (GWAS) meta-analysis identifying 85 LBD risk genes \cite{chia2025gwas}, a 2024 DNA methylation meta-analysis of 1,239 cortical brain samples \cite{nabais2024methylation}, plasma metabolomics distinguishing DLB from Alzheimer's disease (AD) \cite{frontiers2024metabolomics}, and single-cell peripheral immune profiling in LBD \cite{singlecell2024immune}---we propose a systems biology pipeline employing multi-omics factor analysis (MOFA+), similarity network fusion (SNF), and machine learning classifiers to derive cross-layer molecular signatures of LBD. The three primary objectives are: (1) to identify blood-accessible multi-omics biomarker signatures capable of discriminating LBD subtypes from AD and healthy aging; (2) to pinpoint convergent pathogenic pathways across genomic, epigenomic, and metabolic layers; and (3) to prioritize candidate drug targets for future experimental validation. This document provides the scientific rationale, a structured literature review of key recent publications, a detailed proposed methodology, and a feasibility assessment, collectively intended to support an informed decision on whether to initiate this research project.
\end{abstract}

\begin{IEEEkeywords}
Lewy body dementia, dementia with Lewy bodies, multi-omics, data integration, biomarker discovery, alpha-synuclein, GWAS, epigenomics, metabolomics, single-cell immunoprofiling.
\end{IEEEkeywords}
